% Options for packages loaded elsewhere
\PassOptionsToPackage{unicode}{hyperref}
\PassOptionsToPackage{hyphens}{url}
\PassOptionsToPackage{dvipsnames,svgnames,x11names}{xcolor}
\documentclass[
  11pt,
]{article}
\usepackage{xcolor}
\usepackage[margin=27mm]{geometry}
\usepackage{amsmath,amssymb}
\setcounter{secnumdepth}{-\maxdimen} % remove section numbering
\usepackage{iftex}
\ifPDFTeX
  \usepackage[T1]{fontenc}
  \usepackage[utf8]{inputenc}
  \usepackage{textcomp} % provide euro and other symbols
\else % if luatex or xetex
  \usepackage{unicode-math} % this also loads fontspec
  \defaultfontfeatures{Scale=MatchLowercase}
  \defaultfontfeatures[\rmfamily]{Ligatures=TeX,Scale=1}
\fi
\usepackage{lmodern}
\ifPDFTeX\else
  % xetex/luatex font selection
  \setmainfont[]{Cambria}
  \setmonofont[]{Consolas}
\fi
% Use upquote if available, for straight quotes in verbatim environments
\IfFileExists{upquote.sty}{\usepackage{upquote}}{}
\IfFileExists{microtype.sty}{% use microtype if available
  \usepackage[]{microtype}
  \UseMicrotypeSet[protrusion]{basicmath} % disable protrusion for tt fonts
}{}
\makeatletter
\@ifundefined{KOMAClassName}{% if non-KOMA class
  \IfFileExists{parskip.sty}{%
    \usepackage{parskip}
  }{% else
    \setlength{\parindent}{0pt}
    \setlength{\parskip}{6pt plus 2pt minus 1pt}}
}{% if KOMA class
  \KOMAoptions{parskip=half}}
\makeatother
\setlength{\emergencystretch}{3em} % prevent overfull lines
\providecommand{\tightlist}{%
  \setlength{\itemsep}{0pt}\setlength{\parskip}{0pt}}
\usepackage{bookmark}
\IfFileExists{xurl.sty}{\usepackage{xurl}}{} % add URL line breaks if available
\urlstyle{same}
\hypersetup{
  pdftitle={WP-11 · RPP adversarial analysis --- experiment design \& pre-registration (v1)},
  pdfauthor={The Privacy-is-Value Research Programme},
  colorlinks=true,
  linkcolor={black},
  filecolor={Maroon},
  citecolor={Blue},
  urlcolor={blue},
  pdfcreator={LaTeX via pandoc}}

\title{WP-11 · RPP adversarial analysis --- experiment design \&
pre-registration (v1)}
\usepackage{etoolbox}
\makeatletter
\providecommand{\subtitle}[1]{% add subtitle to \maketitle
  \apptocmd{\@title}{\par {\large #1 \par}}{}{}
}
\makeatother
\subtitle{WP-11 (rpp-adversarial-paper) · tier internal (design doc;
feeds the tier-A WP-11 paper) · 2026-07-18 · status DESIGN v1 ---
PRE-REGISTRATION DRAFT, not yet committed}
\author{The Privacy-is-Value Research Programme}
\date{2026-07-18}

\begin{document}
\maketitle

\subsection{0. The iron rule (A6), stated
first}\label{the-iron-rule-a6-stated-first}

Every hypothesis below has a \textbf{pre-registered falsification
condition}: a written, dated statement of the observation that would
make the claim FALSE, fixed before any measurement runs. This design is
willing to refute RPP on every axis, and says so per hypothesis. Three
specific ways this study can fail RPP, named now so they cannot be
explained away later: - \textbf{H1 can fail:} models are trained to
resist injected instructions; if they proceed normally and ignore the
RPP directive, the ``forcing'' claim is false. - \textbf{H2 can fail:} a
malicious counterparty was present at the shared experience by
construction; if their informed guess reconstructs the withheld fraction
at operational σ, the two-party security is false. - \textbf{H3 can
fail:} if relational proverbs are recalled no better than random
mnemonics, the relational-memory claim is false.

Results file to the register whichever way they fall (GR-8). Harness
code must not embed the hoped-for result (A6 failure mode); a second
engineer re-derives the metrics from raw logs.

\subsection{1. What is being tested (the ruled claims, not the old
framing)}\label{what-is-being-tested-the-ruled-claims-not-the-old-framing}

The First Person ruled (L162/L163) that RPP's novelty is: \textbf{M1}
--- prompt injection used as a \emph{constructive} mechanism that forces
a cognitive act as a precondition of use; and the \textbf{Option-C
interleaved bilateral construction} --- the trust object is κ =
hash(interleave(P\_A, P\_B)), the interleaved pair split across two
parties, with a tunable disclosure ratio σ ∈ {[}0,1{]} (the public
anchor carries σ of the interleaved pair; the two parties jointly hold
1−σ); with \textbf{M3} --- relational/episodic memorability of the
proverbs --- as a design rationale. This study measures those three
claims and the σ-frontier that integrates them. It does NOT test
``comprehension as authentication'' (the superseded framing; WP-11a
§1-§2).

Definitions used throughout (GR-4 formal register): the two participants
are \textbf{party A} and \textbf{party B}; each authors a
natural-language compression \textbf{P\_A}, \textbf{P\_B} of a shared
experience; \textbf{interleave(·)} is the third inscription mode
(fragments of P\_A and P\_B interwoven across the visibility boundary);
\textbf{σ} is the fraction of the interleaved pair disclosed in the
public anchor; the \textbf{injected directive} is the RPP instruction
embedded in content requiring a reader to author a binding proverb
before proceeding.

\subsection{2. Hypotheses, metrics, and pre-registered
falsification}\label{hypotheses-metrics-and-pre-registered-falsification}

\subsubsection{H1 --- the forcing claim
(M1)}\label{h1-the-forcing-claim-m1}

\textbf{Hypothesis.} Content carrying the injected RPP directive causes
a reader/model to perform the proverb-formation act (or be detectably
gated from proceeding without it) at a rate materially above a
no-directive control, and this forcing survives a defined
strip/paraphrase attack at a rate materially above a canary baseline.
\textbf{Primary metric.} \emph{Forced-act rate} f = P(model emits a
valid content-bound proverb before answering \textbar{} directive
present) minus the same rate under a no-directive control.
\emph{Strip-survival rate} s = f measured after each of three
adversarial pre-processing attacks (directive deletion, directive
paraphrase, directive relocation). \textbf{Carve-out metric (mandatory,
per WP-11a §0).} \emph{Content-binding rate} b = P(the emitted proverb
is specific to THIS content, failing a generic-completion detector).
This is what distinguishes M1 from self-steering CoT (which forces
reasoning but is not content-bound) and from context-bombs (which derail
rather than compel a cooperative act). A high f with low b would mean
the directive forces \emph{a} proverb but not a \emph{bound} one --- a
partial result stated as such. \textbf{Pre-registered falsification.} M1
is FALSE if, across the model panel, (f ≤ 0.2 at the pre-registered α)
--- models ignore the directive --- OR if directive-deletion alone
reduces f to the canary baseline (the forcing is no stronger than a
passive tripwire), OR if b ≤ 0.5 (the proverb is not content-bound, so
RPP reduces to CoT). \textbf{Panel.} ≥ 6 models across ≥ 3 capability
generations and ≥ 3 vendors, plus an
instruction-tuned-to-resist-injection subset (the hostile case: if
injection-resistant models comply, the constructive-injection claim is
stronger; if only compliant models comply, that is a scope limit stated
honestly).

\subsubsection{H2 --- the two-party interleaved security (Option C, vs
σ)}\label{h2-the-two-party-interleaved-security-option-c-vs-ux3c3}

\textbf{Hypothesis.} A malicious counterparty holding their own half
(P\_B) and the public σ-fraction of the interleaved anchor cannot
reconstruct the withheld (1−σ) fraction of P\_A, nor forge the joint κ,
at better than a pre-registered baseline --- and the reconstruction
difficulty increases as σ decreases. \textbf{Threat model (the honest
one, per L163).} The adversary is a \textbf{malicious counterparty}, not
an outside guesser: they were present at the shared experience, so their
prior over P\_A is \emph{informed}. The benchmark supplies the adversary
the shared-experience context, P\_B, the public σ-fraction, and an
optimisation budget, and measures reconstruction of the withheld
fraction / κ-forgery. \textbf{Primary metric.} \emph{Reconstruction
success} r(σ) = P(adversary recovers the (1−σ) fraction closely enough
to pass the fuzzy-match / recompute κ), measured across σ ∈ \{0.0,
0.382, 0.5, 0.618, 1.0\} (the register C54 candidate values + endpoints,
cited as candidates not optima). The security frontier is r(σ) vs σ.
\textbf{Pre-registered falsification.} The Option-C security is FALSE if
r(σ) exceeds the pre-registered baseline at any operational σ
(specifically if r(0.618) ≥ 0.5 --- a well-informed counterparty
recovers the majority-public sweet-spot), OR if r(σ) does not decrease
monotonically as σ decreases (σ is not a real security knob). A
refutation here is the most consequential possible result and is
reported at contribution prominence. \textbf{Note on scope.} This
measures the SEMANTIC layer only (can the withheld proverb-fraction be
guessed). It does NOT re-test the hard crypto beneath
(SHA-256/Ed25519/κ, whose security is assumed from E11's documented
implementation); the paper states this boundary.

\subsubsection{H3 --- relational recall vs random mnemonics
(M3)}\label{h3-relational-recall-vs-random-mnemonics-m3}

\textbf{Hypothesis.} Human participants regenerate a
relationally-derived proverb closely enough to recover the (1−σ)
fraction, at a rate materially above matched-length random-word
mnemonics, and the gap widens with elapsed time. \textbf{Primary
metric.} \emph{Recall-recovery rate} ρ(t) = P(participant regenerates
their proverb within the fuzzy-match tolerance at delay t ∈ \{1 day, 1
week, 1 month\}), for relational-proverb vs BIP-39-style random-mnemonic
arms of matched entropy/length. \textbf{Pre-registered falsification.}
M3 is FALSE if ρ\_relational(t) ≤ ρ\_random(t) at any t (relational is
not better), OR if ρ\_relational(1 month) is below the recovery
threshold the system needs (memorable but not recoverable-enough to be
usable). This is a usable-security study (SOUPS/USEC form); IRB/consent
required; it shares its design with the WP-11a M3 outstanding sweep.

\subsubsection{The σ-frontier (the integrating
measurement)}\label{the-ux3c3-frontier-the-integrating-measurement}

H2 and H3 both depend on σ and pull opposite ways: higher σ (more
public) → easier recall/verification (H3 up) but weaker security (H2
down); lower σ → stronger security but harder recall. \textbf{The
decisive empirical question is whether a σ exists where H2 holds (r
below baseline) AND H3 holds (ρ above threshold) simultaneously.}
Pre-registered success condition: a non-empty σ-band where both hold.
Pre-registered falsification of the \emph{design} (not just a
hypothesis): if no σ satisfies both, the σ-tunable interleaved
construction does not have an operating point, and the paper reports
that the mechanism is not viable as specified. This is the
``t*-of-comprehension'' crossover the A6 role card anticipated, made
concrete for the ruled construction.

\subsection{3. Conditions / control arms (E3-C34 precedent: multi-arm,
pre-registered)}\label{conditions-control-arms-e3-c34-precedent-multi-arm-pre-registered}

Inheriting E3-C34's seven-control-arm discipline, adapted to the ruled
claims: - \textbf{A0} no-directive control (H1 baseline). - \textbf{A1}
directive present, benign reader (H1 forced-act). - \textbf{A2}
directive present, adversarial strip/paraphrase (H1 strip-survival). -
\textbf{A3} directive present, injection-resistant model (H1 hostile
case). - \textbf{A4} two-party interleave at each σ, honest parties
(H2/H3 baseline). - \textbf{A5} two-party interleave, malicious informed
counterparty (H2 security). - \textbf{A6} relational-proverb recall vs
A7 random-mnemonic recall (H3), each across delays.

\subsection{4. Power, sample, analysis plan
(pre-registered)}\label{power-sample-analysis-plan-pre-registered}

\begin{itemize}
\tightlist
\item
  \textbf{H1:} with 6+ models × 3 attack types × an item corpus of ≥ 200
  content pieces (stratified by domain/length), a two-proportion test
  detects f-gaps ≥ 0.15 at α = 0.01, power 0.9; per-model reported, not
  pooled (models are not exchangeable).
\item
  \textbf{H2:} adversary given a fixed optimisation budget (stated in
  compute-hours and query count); r(σ) estimated over ≥ 100 independent
  (experience, P\_A, P\_B) triples per σ; Clopper-Pearson intervals;
  monotonicity tested by isotonic regression with a pre-registered
  violation threshold.
\item
  \textbf{H3:} between-subjects (relational vs random) × within-subjects
  (delay); target N per arm powered for a 0.15 recall-rate gap at α =
  0.01, power 0.9 (≈ 180/arm pending pilot variance); mixed-effects
  model with participant random effects; attrition plan pre-stated.
\item
  \textbf{Multiplicity:} three hypothesis families, Bonferroni-Holm
  across primary metrics within each; the σ-frontier is a conjunction,
  not a fourth test.
\item
  \textbf{Analysis code} is written against synthetic data with known
  ground truth BEFORE real data exists, and frozen (prevents
  peeking-then-fitting).
\end{itemize}

\subsection{5. Harness design (buildable by a
stranger)}\label{harness-design-buildable-by-a-stranger}

\begin{itemize}
\tightlist
\item
  \textbf{H1 harness:} a corpus loader (content ± injected directive), a
  model-runner (panel via API/local), a \emph{content-binding detector}
  (the generic-completion classifier --- itself validated against a
  labelled set with reported precision/recall so it is not a black box
  embedding the result), and a raw-log store; metrics re-derived from
  logs by a second engineer.
\item
  \textbf{H2 harness:} an interleave(P\_A, P\_B, σ) implementation (the
  design-layer mechanism per E3-C10/C22 --- built here for measurement,
  deployment status stated), a fuzzy-match/κ-recompute verifier, and an
  adversary module (informed prior over P\_A from the shared-experience
  context + an optimisation loop with a fixed budget). The adversary is
  written by someone OTHER than the interleave author (adversarial
  independence).
\item
  \textbf{H3 harness:} a standard usable-security recall protocol
  (enrolment, delayed recall, fuzzy-match scoring), IRB/consent,
  pre-registered on OSF or equivalent.
\item
  All three: seeds fixed and logged; a full run reproducible from a
  committed config; no floating-point non-determinism in the metric
  path.
\end{itemize}

\subsection{6. Deployment-honesty statement (carried into the paper,
GR-8 /
E3-C32)}\label{deployment-honesty-statement-carried-into-the-paper-gr-8-e3-c32}

The interleaved inscription path and the σ disclosure budget are
\textbf{design-layer} in the corpus: the deployed RPP system currently
uses the asymmetric-default inscription and a stratum-derived
visibility, not the σ ∈ {[}0,1{]} budget (E3-C22, E3-C32). This
benchmark BUILDS the interleaved+σ mechanism to measure it; the paper
states plainly that it evaluates a specified-but-not-yet-deployed
construction, and that the φ sweet-spots for σ are register C54
candidate values, not asserted optima. The hard-crypto layer beneath
(SHA-256, Ed25519, κ content-addressing) is assumed from E11's
documented implementation and not re-tested here.

\subsection{7. What this design does NOT do (scope
fences)}\label{what-this-design-does-not-do-scope-fences}

\begin{itemize}
\tightlist
\item
  It does not test the auth-primitive claims (WP-11a §1-§2 superseded).
\item
  It does not re-test the hard cryptographic primitives (H2 measures the
  semantic layer only).
\item
  It does not assert the φ σ-values as optimal (they are candidate
  operating points).
\item
  It does not claim M1 works on all models unconditionally --- the
  hostile injection-resistant arm (A3) exists precisely to bound the
  claim.
\end{itemize}

\subsection{8. Handoff (A6 → next)}\label{handoff-a6-next}

\begin{itemize}
\tightlist
\item
  \textbf{First Person:} commit (date + freeze) this pre-registration
  before any run; the run itself and result-filing are downstream. The
  one design decision still open for you: the \emph{optimisation budget}
  for the H2 adversary (compute-hours + query count) --- this sets how
  hard the malicious-counterparty tries, and it is a policy choice, not
  a technical one. A generous budget makes a passing H2 result stronger
  and a failing one honest.
\item
  \textbf{A3 (formalist, next chain step):} formalise the H2 security as
  a function of the two-party split and (1−σ) --- the benchmark measures
  r(σ); A3 states what bound r(σ) is estimating, and whether the
  interleaved construction has a stated security definition
  (indistinguishability of the withheld fraction given P\_B +
  σ-disclosure + informed prior). H1's ``forcing'' also wants a
  definition A3 can point the metric at.
\item
  \textbf{A4:} obtain TDCommons-9969 full text and verify the three
  memory DOIs (WP-11a §4b); confirm the H2 threat model is genuinely
  distinct from TDCommons's single-narrative derive at the full-claim
  level.
\item
  \textbf{A5-pc:} the adversarial review runs after results, not now;
  but A5 should pre-read this design for the one thing a PC breaks a
  prereg on --- an unfalsifiable hypothesis or a peeking hole. Every H
  here has a written falsification; the σ-frontier can refute the whole
  design.
\end{itemize}

\end{document}
